As future work, we plan to investigate the potential integration of  {\THESYSTEM} in an adaptive data analysis framework like Guess and check by Rogers at al.~\cite{RogersRSSTW20}. As we discussed, this framework is  designed to support adaptive data analyses with limited generalization error. As our experiments show, this framework could benefit from the information provided by {\THESYSTEM} to provide more precise estimate and improved confidence intervals. Another direction we will explore is to make 
the upper bounds provided by {\THESYSTEM} more precise by integrating our algorithm with a path-sensitive approach.

\subsubsection{Accurate Adaptivity Definition}
\label{sec:future-adapt-definition}
The program's adaptivity definition in Section~\ref{sec:adapt-exe}
(specifically in Definition~\ref{def:trace_adapt}), isn't precise enough w.r.t. the intuitive \emph{adaptivity} rounds.
It comes across an over-approximation.
The Example~\ref{ex:multiRoundsS} in Section~\ref{sec:adapt-example} shows this limitation.
% It is 
In this example, the formalized adaptivity by execution-based analysis in Section~\ref{sec:adapt-exe} 
is the initial value of the input variable $k$.
However, the intuitive \emph{adaptivity} is only $2$ given whatever initial trace.
This is resulted from the
%  difference 
irrelevance between the 
% dependency Depth 
dependency quantity (i.e., the vertex weight)  and 
the data dependency relation (the edge) in \emph{variable may-dependency} definition in Section~\ref{sec:dynamic-datadep}.
Specifically the irrelevance between the vertex weight and the edge on the \emph{semantics-based dependency} graph.
It occurs when the 
% weight 
dependency quantity is computed over the traces different from the traces used in 
% witness 
analyzing the \emph{variable may-dependency} relation.

\paragraph{Proposed Methodology}
\label{para:furthers-dep-depth}
Based on the \emph{adaptivity} formalization in Section~\ref{sec:adapt-exe}, the improvement on the accuracy 
of
\emph{adaptivity} definition is planned to develop in following three-steps.
\begin{enumerate}
\item In the first stage of formalization, 
I will give an alternative variable \emph{may-dependency} definition 
by referring to the analysis methodology in \cite{Cousot19a}.
%
Specifically, I will define the variables' dependency relation over two witness traces and an initial trace. Comparing to 
the existing \emph{may-dependency} definition, which quantifies overall possible execution traces, the alternative
the definition explicitly relies on two specific witness traces from the program executions.
This externalization helps in analyzing the \emph{dependency quantity} based on the same 
witness traces as the \emph{may-dependency} relation. In this way, the over-approximation as illustrated above
can be reduced.
%
\item Then based on the new \emph{may-dependency} definition,
I will add the weight of every edge constructed from 
\emph{may-dependency} relation in the \emph{semantics-based dependency} graph, w.r.t. to the witness traces.
%
\item Then in the third stage, the \emph{adaptivity} is still defined as the 
length of the longest finite walk. Differently, the finite walk is added
with extra restriction.
The occurrence of every edge in a finite walk is restricted to be no more than its weight as well.
\end{enumerate}
Through the three steps above, ideally we will give a more accurate definition of the intuitive \emph{adaptivity}.

\subsubsection{Accurate Adaptivity Estimation}
\label{sec:future-adaptfun}
In static program analysis framework $\THESYSTEM$, specifically on the dependency quantity, 
I adopt the reachability bound analysis technique to estimate this dependency quantity.
% In existing static reachability bound analysis, 
However, it isn't precise enough w.r.t. the execution-based reachability bound on every program command.
It comes across an over-approximation on estimation due to its path-insensitive nature.
The Example~\ref{ex:multiRoundsO} shows this limitation in Section~\ref{sec:adapt-example}.
% It is 
In this example, the estimated adaptivity by static analysis in Section~\ref{sec:adapt-static} 
is $1 + 2*k$ where $k$ is the initial value of the input variable.
However, the formal \emph{adaptivity} is only $1 + k$ from the execution-based analysis.

This is resulted from the  path-insensitive of the dependency quantity in
second stage of the static analysis in Section~\ref{sec:alg_weightgen}.
It occurs when the control flow can be decided in a particular way in front of conditional branches, 
while the static analysis fails to witness. 

\paragraph{Proposed Methodology}
\label{para:future-adaptfun}
The accurate static adaptivity analysis is planned to be improved in following two steps.
\begin{enumerate}
    \item \textbf{Path-sensitive Reachability-bound Algorithm Design}
    The imprecision comes from the second stage of the static program analysis.
    The algorithm in this stage is to statically estimate the 
    reachability bound for every location of program.
    Specifically, this reachability bound analysis algorithm isn't path sensitive. 
    There are many works in the program complexity analysis area, estimating the path-sensitive loop bound 
    or reachability bound
    \cite{GustafssonEL05, HumenbergerJK18}, 
    \cite{BrockschmidtEFFG16,AlbertAGP08,AliasDFG10,Flores-MontoyaH14}, 
    \cite{GulwaniZ10, SinnZV17,GulwaniJK09, GulwaniMC09, abs-2203-04243}. 
    However, they aren't analyzing the reachability
    bound in a path-sensitive way.
    \\
    Motivated by this observation, I will first design path-sensitive reachability bound analysis algorithm computing the 
    reachability bounds for every labeled command taking the different paths inside while loop into consideration.
    \item \textbf{$\THESYSTEM$ via Path-sensitive Reachability-bound Algorithm Design}
    Then following the same static program analysis based $\THESYSTEM$ framework in Section~\ref{sec:adapt-static},
   this improved analysis result will be applied to estimate the dependency quantity as follows,
\begin{itemize}
    \item $\THESYSTEM$ will construct different abstract control flow graph based on the requirement of the newly-designed
    path-sensitive reachability bound algorithm.

    \item The improved $\THESYSTEM$ will perform the same algorithm for analyzing the dependency relation in the first stage.
    \item But in the second stage, the improved $\THESYSTEM$ 
    apply the newly-designed  path-sensitive reachability bound algorithm for analyzing the dependency quantity.
    The result from the new algorithm 
    give a tighter approximation on the dependency quantity,
    specifically on the reachability times of the program's every location.
    \item Based on the two results analyzed from the steps above, 
    i.e., the estimated dependency relation and the tighter reachability bound, the improved $\THESYSTEM$ will construct
a more accurate program-based data dependency graph.
Then, by applying Algorithm~\ref{alg:adapt} in Section~\ref{sec:alg_adaptcompute}, 
the improved $\THESYSTEM$ will compute a tighter \emph{adaptivity} for a program.%
\end{itemize}
\end{enumerate}
Again as shown in Figure~\ref{fig:multiRoundsO}(b), ideally we will compute $\frac{k}{2}$ as the weight for vertices $p^6$ and $x^7$  through the newly-designed path-sensitive reachability bound analysis algorithm!
This also gives the motivation for the work in \redd{ PART \romannum{2}} in this thesis.

\subsubsection{Adaptivity Analysis Practical Application}
\label{sec:future-adapt-app}
